% PAGES: 133

% Continuation of sec04_c1_3.tex (page 132 / folio 116): part (iii) below
% continues the manual "Examples:" block opened there (diagonal / lower
% triangular / upper triangular matrices). Do not print a new "Examples:"
% label.
(iii) Let $U$ be an upper triangular matrix of size $n$. The characteristic polynomial
of $U$ is
\begin{equation}
P_U(t) \;=\; \det(tI-U) \;=\; \prod_{i=1}^n (t-U(i,i));
\end{equation}
cf. (10.4). Since the roots of $P_U(t)$ are $U(i,i)$, $i = 1:n$, we conclude that the
eigenvalues of $U$ are its diagonal entries $U(i,i)$, $i = 1:n$. Note that the
eigenvector corresponding to the eigenvalue $U(1,1)$ is $e_1$, the first column of the
identity matrix. However, the eigenvector corresponding to the eigenvalue $U(i,i)$, with
$2 \le i \le n$, is not necessarily $e_i$, the $i$-th column of the identity matrix.
\hfill$\square$\par\medskip

\begin{definition}
If $\lambda$ is a root of multiplicity $m_\lambda$ of the characteristic polynomial
$P_A(t)$ corresponding to matrix $A$, then, by definition, $m_\lambda$ is the
multiplicity of the eigenvalue $\lambda$ of $A$.
\end{definition}

\begin{theorem}
A square matrix of size $n$ has $n$ eigenvalues, counted with their multiplicities;
some of these eigenvalues may be complex numbers.

In other words, if $A$ is a square matrix of size $n$, then
\[
\sum_{\lambda \in \lambda(A)} m_\lambda \;=\; n,
\]
where $\lambda(A)$ is the set of all the eigenvalues\footnote{The set of all the
eigenvalues of a matrix is also called the spectrum of the matrix.} of $A$.
\end{theorem}

As a direct consequence of Theorem 4.2, note that a square matrix of size $n$ has at
most $n$ distinct eigenvalues.

\begin{proof}
If $A$ is a square matrix of size $n$, then its characteristic polynomial $P_A(t)$ has
degree $n$. The roots of $P_A(t)$ are the eigenvalues of $A$; cf. Theorem 4.1. Recall
from the Fundamental Theorem of Algebra that a polynomial of degree $n$ with real
coefficients has exactly $n$ roots when counted with their multiplicities. Note that the
roots of the polynomial may be complex numbers.

From Definition 4.3 of the multiplicity of an eigenvalue, and since the sum of the
multiplicities of the roots of a polynomial of degree $n$ is equal to $n$, we conclude
that the number of the eigenvalues of $A$, counted with their multiplicities, is $n$.
\end{proof}

\begin{definition}
The number of different eigenvectors corresponding to an eigenvalue $\lambda$ is equal
to the largest number of linearly independent eigenvectors corresponding to $\lambda$.
In other words, if $V_\lambda = \{v \ \text{such that} \ Av = \lambda v\}$ is the space
of the eigenvectors of the matrix $A$ corresponding to $\lambda$, then the number of
different eigenvectors corresponding to an eigenvalue $\lambda$ is equal to
$\dim(V_\lambda)$, where $V_\lambda$ denotes the dimension of the space $V_\lambda$; see
Definition 1.11.
\end{definition}

Every eigenvalue has at least one corresponding eigenvector. If an eigenvalue $\lambda$
has multiplicity $m_\lambda$, then there could exist at most $m_\lambda$ linearly
independent eigenvectors corresponding to $\lambda$. However, it might also happen that
there are less than $m_\lambda$ eigenvectors corresponding to $\lambda$, and therefore a
matrix might have less than $n$ linearly independent eigenvectors; see the example at
the beginning of this section, as well as the example below:
% PDF page 134 (not part of this assignment) continues here with an
% unnumbered worked example -- no environment is opened for it in this
% file; the sentence above ends in a colon leading into that example.
